\section{Fourier transform formulas}\label{appendixC}

We f\/ix complex numbers $\mu,\nu$ satisfying
\begin{gather}\label{munu}
-a<\im \mu<\im \nu <a,\qquad a=(a_{+}+a_{-})/2.
\end{gather}
Hence we have
\begin{gather}
\im (\nu-\mu)\in(0,a_{+}+a_{-}),
\end{gather}
and the function
\begin{gather}\label{defI}
I(\mu,\nu;x)\equiv G(x-\nu)/G(x-\mu),
\end{gather}
is pole-free in the strip
\begin{gather}\label{xstrip}
\im \nu -a<\im x<\im \mu +a.
\end{gather}
Also, from the $G$-asymptotics~\eqref{Gas} we deduce
\begin{gather}
I(\mu,\nu;x)=O(\exp(\mp \alpha\im (\nu -\mu)\re x/2)),\qquad \re x\to\pm\infty.
\end{gather}
Therefore, for real $y$ the function
\begin{gather}\label{defFy}
F(\mu,\nu;y)\equiv \int_{\R}dx \exp(i\alpha xy)\frac{G(x-\nu)}{G(x-\mu)},
\end{gather}
is well def\/ined, and analytic in $\mu$ and $\nu$ in the region~\eqref{munu}. Moreover, we retain exponential decay of the integrand when we let $y$ vary over the strip
\begin{gather}\label{ystrip}
|\im y|<\im(\nu-\mu)/2,
\end{gather}
so $F(\mu,\nu;y)$ is analytic in $y$ in this strip and given by~\eqref{defFy}.

Our f\/irst aim is to obtain the Fourier transform~\eqref{defFy} in explicit form. Since the func\-tion~$I$~\eqref{defI} has no poles in the strip~\eqref{xstrip}, we can make a contour shift
\begin{gather}
x \to z +(\mu+\nu)/2,\qquad x,z\in\R,
\end{gather}
to deduce
\begin{gather}\label{Fkap}
F(\mu,\nu;y)=\exp(i\alpha y(\mu+\nu)/2)\int_{\R}dz \exp(i\alpha zy)G(\pm z-\kappa),
\end{gather}
where we have set
\begin{gather}\label{kappa}
\kappa \equiv (\nu-\mu)/2,\qquad \im \kappa \in(0,a).
\end{gather}
Hence calculating $F$ amounts to f\/inding the cosine transform of $G(\pm z-\kappa)$ for $\kappa$ in the strip $\im \kappa\in(0,a)$. The reasoning in the proof of the following proposition, however, hinges on staying at f\/irst with~\eqref{defFy}.

A special case of the following result amounts to a hyperbolic analog of a trigonometric beta integral in Ramanujan's lost notebook. More precisely, \eqref{FFt} with $y=(\mu+\nu)/2$ amounts to equation~(1.8) with $\tau<0$ in Stokman's paper~\cite{stok}. Formula~\eqref{FFt} is also obtained in~Chapter~5 of van de Bult's Ph.D. thesis~\cite{bult2} (cf.~Theorem~5.6.8 with $n=1$), as a result of specializing more general formulas. In slightly dif\/ferent guises it occurred previously in various papers (the earliest ones being~\cite{faka,kash,pote}), but a complete proof cannot be found there.

\begin{prop}\label{propositionC.1}
The Fourier transform~\eqref{defFy} admits the explicit evaluation
\begin{gather}\label{FFt}
F(\mu,\nu;y)=(a_{+}a_{-})^{1/2}\exp(i\alpha y(\mu+\nu)/2)G(ia-2\kappa)G(\pm y-ia+\kappa),
\end{gather}
where $\mu,\nu$ satisfy~\eqref{munu}, $y$ satisfies \eqref{ystrip}, and $\kappa$ is given by~\eqref{kappa}.
\end{prop}
\begin{proof}
We f\/irst take $y$ real and choose $\mu,\nu$ such that
\begin{gather}\label{ass}
\im(\nu-\mu)\in  (a_l,a_{+}+a_{-}),\qquad a_l\equiv \max(a_{+},a_{-}).
\end{gather}
This entails that we can shift $\nu$ down by $ia_{\de}$ or $\mu$ up by $ia_{\de}$ without leaving the region~\eqref{munu}. Hence we may shift under the integral sign and use the $G$-A$\De$Es~\eqref{Gades} to obtain
\begin{gather}
F(\mu,\nu-ia_{\de};y)=2\int_{\R}dxe^{i\alpha xy}c_{-\de}(x-\nu+ia_{\de}/2)G(x-\nu)/G(x-\mu),\qquad y\in\R.
\end{gather}
Recalling that $F(\mu,\nu;y)$ is analytic in $y$ when~\eqref{ystrip} holds, this can be rewritten as
\begin{gather}
F(\mu,\nu-ia_{\de};y)=e_{-\de}(-\nu+ia_{\de}/2)F(\mu,\nu;y-ia_{\de}/2)\nonumber\\
\phantom{F(\mu,\nu-ia_{\de};y)=}{}
+e_{-\de}(\nu-ia_{\de}/2)F(\mu,\nu;y+ia_{\de}/2).\label{Fs1}
\end{gather}
Likewise, we obtain
\begin{gather}
F(\mu+ia_{\de},\nu;y)=e_{-\de}(-\mu-ia_{\de}/2)F(\mu,\nu;y-ia_{\de}/2)\nonumber\\
\phantom{F(\mu+ia_{\de},\nu;y)=}{}
+e_{-\de}(\mu+ia_{\de}/2)F(\mu,\nu;y+ia_{\de}/2).\label{Fs2}
\end{gather}

Next, consider the contour shift $x\to x+is$, where
\begin{gather}
s\in (\im\nu-a,\im \mu+a),
\end{gather}
cf.~\eqref{xstrip}. It implies
\begin{gather}
F(\mu,\nu;y)     =     \int_{\R}dx\exp(i\alpha (x+is)y)G(x+is-\nu)/G(x+is-\mu)
\nonumber \\
\phantom{F(\mu,\nu;y)     =}{}  =    \exp(-\alpha sy)F(\mu-is,\nu-is;y),\qquad y\in\R.\label{xshi}
\end{gather}
Since $F(\mu,\nu;y)$ is analytic in $y$ in the strip~\eqref{ystrip}, it now follows from~\eqref{ass} that we have
\begin{gather}\label{Fkey}
F(\mu-is,\nu-is;y)=e^{\alpha sy}F(\mu,\nu;y),\qquad |\im y|\le a_l/2.
\end{gather}

Now in~\eqref{Fs2} we may take $\mu,\nu\to \mu-is,\nu-is$ with $s\in(0,\im\mu +a)$. Then we are entitled to invoke~\eqref{Fkey} to get
\begin{gather}
F(\mu-is+ia_{\de},\nu-is;y)     =     e_{-\de}(-\mu+is-3ia_{\de}/2)e^{\alpha sy}F(\mu,\nu;y-ia_{\de}/2)
\nonumber \\
\phantom{F(\mu-is+ia_{\de},\nu-is;y)     =}{}   +e_{-\de}(\mu-is+3ia_{\de}/2)e^{\alpha sy}F(\mu,\nu;y+ia_{\de}/2).
\end{gather}
In this equation we can let $s$ converge to $a_{\de}$, which yields
\begin{gather}
F(\mu,\nu-ia_{\de};y)     =     e_{-\de}(2y-\mu-ia_{\de}/2)F(\mu,\nu;y-ia_{\de}/2)
\nonumber \\
\phantom{F(\mu,\nu-ia_{\de};y)     =}{}  +e_{-\de}(2y+\mu+ia_{\de}/2)F(\mu,\nu;y+ia_{\de}/2).\label{Fs3}
\end{gather}

Comparing~\eqref{Fs3} and~\eqref{Fs1}, we see that the dif\/ference yields a linear relation between $F(\mu,\nu;y+ia_{\de}/2)$ and $F(\mu,\nu;y-ia_{\de}/2)$. After some simplif\/ication, this relation can be rewritten as
\begin{gather}\label{Fade}
\frac{F(\mu,\nu;y+ia_{\de}/2)}{F(\mu,\nu;y-ia_{\de}/2)}=e_{-\de}(-\mu-\nu)\frac{c_{-\de}(y-ia+\kappa)}{c_{-\de}(y+ia-\kappa)}.
\end{gather}
Introducing
\begin{gather}\label{Fr}
F_r(\mu,\nu;y)\equiv \exp(i\alpha (\mu+\nu)y/2)G(\pm y-ia+\kappa),
\end{gather}
it is easy to verify that $F_r$ also satisf\/ies this $y$-A$\De$E. Since we can choose $\de=+,-$ and $a_{+}/a_{-}\notin\Q$, it readily follows that we must have
\begin{gather}\label{Fform}
F(\mu,\nu;y)=C(\mu,\nu)F_r(\mu,\nu;y),
\end{gather}
with $C$ independent of $y$.

The upshot of our reasoning thus far is that when $\mu$ and $\nu$ are restricted by~\eqref{ass}, then the Fourier transform is of the form~\eqref{Fform}, with $F_r$ given by~\eqref{Fr}. By analyticity in $\mu$ and $\nu$, this relation now extends to the whole region~\eqref{munu} and then, by analyticity in $y$, to the strip~\eqref{ystrip}. In view of~\eqref{Fkap}, we therefore have shown
\begin{gather}
\int_{\R}dz\exp(i\alpha zy)G(\pm z-\kappa)=C(\mu,\nu)G(\pm y-ia+\kappa),\\
 \im \kappa \in (0,a),\qquad |\im y|<\im \kappa.\nonumber
\end{gather}
It follows from this that we have
\begin{gather}
C(\mu,\nu)=D(\mu-\nu),
\end{gather}
so it remains to prove
\begin{gather}
D(\mu-\nu)=(a_{+}a_{-})^{1/2} G(ia+\mu-\nu).
\end{gather}

To this end we substitute
\begin{gather}
F(\mu,\nu;y)=D(\mu-\nu)\exp(i\alpha (\mu+\nu)y/2)G(\pm y-ia+\kappa)
\end{gather}
in~\eqref{Fs1}, and divide the result by
\begin{gather}
D(\mu-\nu)\exp(i\alpha (\mu+\nu)y/2)G(\pm y-ia_{\de}/2-ia+\kappa).
\end{gather}
Using the $G$-A$\De$Es, a straightforward calculation then yields
\begin{gather}\label{Dade}
\frac{D(\mu-\nu+ia_{\de})}{D(\mu-\nu)}=2is_{-\de}(\mu-\nu+ia_{\de}).
\end{gather}
Setting
\begin{gather}
D_r(\mu-\nu)\equiv G(\mu-\nu+ia),
\end{gather}
we see that $D_r$ also satisf\/ies the A$\De$E~\eqref{Dade}. Thus we deduce as before
\begin{gather}
D(\mu-\nu)=\eta D_r(\mu-\nu),
\end{gather}
where $\eta$ can only depend on $a_{+}$ and $a_{-}$.

As a result, we have now proved
\begin{gather}\label{Ftr}
\int_{\R}dx\exp(i\alpha xy)G(\pm x-\kappa)=\eta G(ia-2\kappa)G(\pm y-ia+\kappa),\\
 \im \kappa \in (0,a),\qquad |\im y|<\im \kappa.\nonumber
\end{gather}
To calculate $\eta$, we choose $\kappa$ equal to $ia_{-}/2$. Using the $G$-A$\De$Es, this yields
\begin{gather}
\int_{\R}dx\exp(i\alpha xy)\frac{1}{c_{+}(x)}=\eta G(ia_{+}/2-ia_{-}/2)\frac{1}{c_{-}(y)}.
\end{gather}
Now the integral is elementary, yielding $a_{+}/c_{-}(y)$. Using~\eqref{Geval} we then infer $\eta$ equals $(a_{+}a_{-})^{1/2}$, hence completing the proof.
\end{proof}
